Climate misinformation can alter the information environment in which investors and liquidity providers trade. This paper examines that relationship in a panel of 204 listed climate-transition securities, using sector-specific exposure constructed from 217,954 English-language Facebook climate posts classified against nine energy sub-industries. Greater exposure is positively associated with quoted bid--ask spreads under monthly and local market-day controls. On a common sample of 486,655 security-days, the market-day estimate is $+0.0129$ percentage points per index point ($p=0.033$ under sector/date clustering), a pattern consistent with information-verification costs. The association changes sign and is imprecise after sector-month conditions are absorbed; with prior-session financial controls, the estimate is $-0.0014$ ($p=0.533$). Timing and alternative-outcome checks do not establish an additional robust response. The findings connect the political production of climate misinformation to the trading environment of energy equities while locating the positive spread association in shared sector conditions. The contribution is a sector-sensitive measure and a transparent account of where the liquidity association is observed, rather than a causal estimate of misinformation costs.
